\section{Introduction}
Underground Hydrogen Storage (UHS) in depleted gas reservoirs and deep saline aquifers has emerged as a cornerstone technology for seasonal green energy storage, enabling the decoupling of intermittent renewable electricity generation from fluctuating market demands~\cite{verma2024uhs}. However, the physical and operational realities of UHS depart substantially from conventional natural gas storage~\cite{pal2024uhsLithuania}. The low viscosity and density of hydrogen relative to resident fluids promote severe gravitational segregation and viscous fingering instabilities, while geochemical and microbiological interactions introduce potential loss mechanisms~\cite{heinemann2021}. Furthermore, balancing seasonal supply demands requires frequent, highly transient injection-withdrawal cycles. The resulting flow fields are characterized by rapid pressure and saturation fluctuations that propagate non-linearly through the reservoir porous medium~\cite{rashid2023geoenergy}.

To model and predict these dynamics, reservoir engineers rely on numerical simulators (e.g., ECLIPSE, INTERSECT, or tNavigator) that solve coupled systems of non-linear partial differential equations representing mass conservation and multi-phase flow via the Finite Volume Method (FVM). Although high-fidelity simulators provide the necessary mathematical rigor, their computational overhead remains prohibitive for multi-query applications. Bayesian history matching, closed-loop schedule optimization, and uncertainty quantification studies demand ensembles of $\mathcal{O}(10^2$--$10^4)$ forward simulations. Executing these ensembles under conventional simulators can require days or weeks of wall-clock time, creating an acute demand for fast, physically consistent surrogate emulators~\cite{pal2023ccs, pal2021mlhistory}.

Scientific Machine Learning (SciML) offers a promising route for accelerating these workflows~\cite{farchi2021}. Convolutional Neural Networks (CNNs) and operator-learning frameworks like the Fourier Neural Operator (FNO) have demonstrated significant speedups for structured, Cartesian grids. However, geological reservoirs are rarely discretized on uniform grids. Perpendicular Bisection (PEBI) and Voronoi tessellations are the industry standard for geological modeling. These unstructured meshes conform naturally to complex geological faults, pinch-outs, and deviated well trajectories, providing local refinement around wellbores without global cell expansion, and satisfying the orthogonality requirements of two-point flux approximation (TPFA) schemes~\cite{ahmed2015mpfa}. Applying grid-based neural architectures to PEBI meshes requires spatial interpolation onto structured grids. This interpolation step introduces severe truncation errors, particularly near wellbores where the pressure gradients are steepest and physical accuracy is most critical.

Graph Neural Networks (GNNs) provide a native solution to this geometric challenge~\cite{geometric}. By representing the irregular grid as a graph $G = (V, E)$ where grid cells map to nodes and shared cell faces map to edges, GNNs can operate directly on the native reservoir mesh. Edge features derived from physical transmissibilities can encode the conductivity of fluid pathways, mapping the physical connectivity in a form directly analogous to the finite-volume stencil~\cite{pal2006tensor}. Thus, GNN message passing can be framed as a learned generalization of the FVM flux update.

This study addresses four interrelated bottlenecks in the development of GNN-based surrogates for UHS reservoirs on PEBI meshes:
\begin{enumerate}
    \item \textbf{The Autocorrelation Trap}: We demonstrate that direct next-step absolute pressure prediction ($P_{t+1}$) is a deceptive target because temporal autocorrelation allows trivial, physics-free models to achieve $R^2 > 0.99$. We quantify the performance gap and advocate for the pressure-change target ($\Delta P$).
    \item \textbf{Temporal Inertia Analysis}: Through systematic feature engineering, we reveal that local temporal lag features (such as $\Delta P_{t-1}$) dominate predictive performance, and we provide a physical explanation in terms of transient pressure diffusion.
    \item \textbf{Native Graph Representation}: We expose how structured-grid spatial features silently fail on PEBI meshes due to integer factorization assumptions and show that graph representations natively resolve this mismatch by incorporating physical transmissibilities as edge weights.
    \item \textbf{Autoregressive Rollout Optimization}: We evaluate Spatio-Temporal GNNs in a 142-step autoregressive rollout mode. We demonstrate that residual skip connections, Layer Normalization, and scheduled sampling curriculum stabilize long-term predictions, yielding a 31.6\% reduction in pressure RMSE.
\end{enumerate}
The manuscript is organized as follows: Section 2 outlines the governing equations and PEBI discretization; Section 3 describes the dataset; Section 4 analyzes the baseline models and the autocorrelation trap; Section 5 formulates the pressure-change target; Section 6 details feature engineering; Section 7 outlines the graph representation; Section 8 describes the ST-GNN architecture; Section 9 presents rollout results; Section 10 discusses the findings; Section 11 presents practical implications; Section 12 details limitations; and Section 13 concludes the work.
